% Project architecture report (pdflatex-friendly)
% This version avoids fontspec/polyglossia/minted so it can be built with pdflatex.
\documentclass[12pt,a4paper]{article}

\usepackage[margin=2.2cm]{geometry}
\usepackage[T2A]{fontenc}
\usepackage[utf8]{inputenc}
\usepackage[russian,english]{babel}

\usepackage{microtype}
\usepackage{hyperref}
\usepackage{bookmark}
\usepackage{graphicx}
\usepackage{longtable}
\usepackage{array}
\usepackage{booktabs}
\usepackage{enumitem}
\usepackage{titlesec}
\usepackage{fancyhdr}
\usepackage{xcolor}
\usepackage{tikz}
\usetikzlibrary{arrows.meta,positioning,fit}

\hypersetup{
  colorlinks=true,
  linkcolor=blue!55!black,
  urlcolor=blue!55!black,
  citecolor=blue!55!black,
  pdftitle={Отчёт по проекту avtodiler (avtodiler-3)},
}

\setlength{\parindent}{0pt}
\setlength{\parskip}{6pt}

\pagestyle{fancy}
\fancyhf{}
\lhead{avtodiler-3}
\rhead{Архитектурный отчёт (pdflatex)}
\cfoot{\thepage}

\title{\textbf{Отчёт по проекту \texttt{avtodiler-3}}\\[4pt]
\large Pdflatex-совместимая версия (без minted/fontspec)}
\author{Репозиторий: \texttt{Sultan-IT-Solutions/avtodiler}}
\date{\today}

\begin{document}
\maketitle

\tableofcontents
\newpage

\section{Краткое описание}
Проект \texttt{avtodiler-3} --- SPA (React/Vite) + API на Vercel (единая serverless функция) + Neon Postgres (JSONB).
Ключевое требование: \textbf{никакого localStorage}, все данные читаются/пишутся в БД через API.

\section{Стек}
\begin{itemize}[leftmargin=18pt]
  \item Frontend: React 18, TypeScript, Vite, React Router, TailwindCSS, i18next.
  \item Backend: Vercel Serverless Function (Node.js ESM), роутер \texttt{api/index.ts}.
  \item DB: Neon Postgres, таблицы \texttt{cars/offers/dealers/leads} с \texttt{data jsonb}.
\end{itemize}

\section{Ограничение Vercel и единая функция}
В \texttt{vercel.json} настроены rewrites, чтобы \texttt{/api/*} попадало в один entrypoint \texttt{/api}.

\begin{verbatim}
{
  "rewrites": [
    { "source": "/api/(.*)", "destination": "/api" },
    { "source": "/(.*)", "destination": "/index.html" }
  ]
}
\end{verbatim}

\section{Сессия админа без localStorage}
Авторизация реализована через HttpOnly cookie \texttt{admin\_session} (подписанный токен). Проверка сессии: \texttt{GET /api/admin/login/status}.

\section{Как собрать PDF}
\begin{verbatim}
# pdflatex
pdflatex -output-directory=docs docs/architecture-report-pdflatex.tex
pdflatex -output-directory=docs docs/architecture-report-pdflatex.tex
\end{verbatim}

\textit{Примечание:} эта версия нарочно короче и без подсветки кода. Полная версия отчёта: \texttt{docs/architecture-report.tex} (XeLaTeX).

\end{document}
